\section{Latency: the lag MPPI does not see}
\label{sec:latency}

\subsection{The model}
Between a command and the odometry that shows its effect sit three delays in series: the pure
delay of the command chain $L_d$ (pulse out, ESC), the thruster spool (first order, time
constant $\tau_s$) and the odometry latency $L_o$ (EKF and transport); then the hull itself
responds with time constant $\tau_h$ ($I_z$ over the linearized damping, for yaw). Seen from
the odometry, a yaw onset from rest is close to a first-order-plus-dead-time (FOPDT)
response,
\[
  r(t) = r_0 + K\Big(1 - e^{-(t - t_\text{on} - L)/\tau}\Big)\quad (t > t_\text{on} + L),
  \qquad r(t) = r_0 \text{ before},
\]
with the apparent dead time
\[
  L \approx L_d + L_o + c\,\tau_s, \qquad c \approx \tfrac12,
\]
since a FOPDT fit of a two-lag cascade moves about half the smaller lag into the dead time
(Skogestad's half rule). On synthetic onsets the fit recovered $L_d + L_o$ within $\pm0.01$\,s,
with $c$ between 0.40 and 0.50 \ev{latency synthetic, 16 cases}. That $L$ is what the loop
feels, so it is the right input to the rules; the closed-loop velocity lag $\tau$ per axis is
measured separately by Controller steps as the 10--90\% rise time over $\ln 9$ (the rise time
of a first-order response is $\tau\ln 9$). The total lag of axis $a$ is
\[
  T_a = \tau_a + L .
\]

\subsection{The fit}
Onsets are pulse outputs that leave neutral ($|p - 1500| \ge 25\,\mu$s on any thruster) after at
least 1\,s of rest. For each onset the yaw rate over $[-0.5, 3]$\,s is fitted on a grid of
$(L, \tau)$, $L \in [0, 0.8]$\,s by 0.01, $\tau \in [0.04, 1.5]$\,s by 0.02, with $r_0$ and $K$ solved
by least squares at each node. An onset says nothing, and is dropped, when $|K|$ is under four
times the pre-onset noise or the optimum sits on the grid's edge. The result is the median $L$
over onsets with a bootstrap 90\% interval (2000 resamples, seed 0).

Yaw is used because its $\tau_h$ is the shortest (0.1--0.3\,s), which separates it best from $L$,
and the gyro is the cleanest channel. Both series are bag receipt times on the boat computer,
so no clock alignment enters.

\paragraph{Acceptance.} At least 6 fitted onsets and a 90\% interval no wider than 0.05\,s;
otherwise one more Characterize yaw pair (it runs both signs).

\paragraph{Evidence in simulation.} In crane\_ml, with \code{pwm\_relay.py} adding a known dead
time (8 onsets each): $L = 0.165$\,s $[0.16, 0.175]$ at relay dead time 0, and $0.29$\,s at 0.1\,s,
so $\Delta L = 0.125$\,s; the relay's 50\,Hz tick accounts for the 0.025\,s excess
\ev{crane relay latency fit}. The sim's own 0.165\,s supports the rule constant
\code{nav2\_rules.delay\_s} $= 0.15$. The 0.2\,s case is not yet run; the validation asks for
$\Delta L$ within $\pm0.03$\,s at 0.1 and 0.2\,s.

\paragraph{When the latency is the problem.} Dead time is the one difference nothing downstream
absorbs (Table~\ref{tab:classify}): 0.3\,s of it raises the yaw step overshoot from 10\% to 82\%. Above
$L = 0.3$\,s the yaw rule of Section~\ref{sec:limits} drops $w_z$ under 0.28\,rad/s at the precise
tolerance, so the procedure stops and fixes the latency (odometry rate, the EKF chain) instead
of tuning around it.

\begin{figure}[tbp]
\centering
\begin{tikzpicture}
\begin{groupplot}[group style={group size=2 by 1, horizontal sep=1.7cm},
  width=0.47\linewidth, height=5.2cm, xmin=-0.015, xmax=0.36, ymin=0,
  xlabel={added delay (s)}, xtick={0,0.1,0.2,0.3},
  tick label style={font=\footnotesize}, label style={font=\small},
  legend cell align=left,
  legend style={font=\footnotesize, draw=none, fill=white, fill opacity=0.85, text opacity=1,
    at={(0.03,0.97)}, anchor=north west},
  grid=major, grid style={gray!20}, clip=false]
\nextgroupplot[ylabel={yaw step overshoot (\%)}, ymax=100]
  \fill[sD!18] (axis cs:0.3,0) rectangle (axis cs:0.36,100);
  \draw[sD, dashed] (axis cs:0.165,0) -- (axis cs:0.165,100);
  \addplot[sA, thick, mark=*, error bars/.cd, y dir=both, y explicit]
    table[x=delay, y=val, y error minus=lo, y error plus=hi] {figures/tuning/data/delay_overshoot_dead.dat};
  \addlegendentry{dead time}
  \addplot[sB, thick, dashed, mark=square*, error bars/.cd, y dir=both, y explicit]
    table[x=delay, y=val, y error minus=lo, y error plus=hi] {figures/tuning/data/delay_overshoot_latency.dat};
  \addlegendentry{odometry latency}
  \node[font=\scriptsize, sD!30!black, anchor=south] at (axis cs:0.165,100) {crane $L \approx 0.165$\,s};
  \node[font=\scriptsize, rotate=90, anchor=north east] at (axis cs:0.3,98) {fix latency first};
\nextgroupplot[ylabel={route yaw reversals (/min)}, ymax=11]
  \fill[sD!18] (axis cs:0.3,0) rectangle (axis cs:0.36,11);
  \draw[sD, dashed] (axis cs:0.165,0) -- (axis cs:0.165,11);
  \addplot[sA, thick, mark=*, error bars/.cd, y dir=both, y explicit]
    table[x=delay, y=val, y error minus=lo, y error plus=hi] {figures/tuning/data/delay_yawrev_dead.dat};
  \addplot[sB, thick, dashed, mark=square*, error bars/.cd, y dir=both, y explicit]
    table[x=delay, y=val, y error minus=lo, y error plus=hi] {figures/tuning/data/delay_yawrev_latency.dat};
  \node[font=\scriptsize, sD!30!black, anchor=south] at (axis cs:0.165,11) {crane $L \approx 0.165$\,s};
  \node[font=\scriptsize, rotate=90, anchor=north east] at (axis cs:0.3,10.8) {fix latency first};
\end{groupplot}
\end{tikzpicture}
\caption{Delay is what the velocity loop cannot absorb. Left: yaw step overshoot against
added dead time (blue) or odometry latency (orange), plant\_v1 seed~0: 9.8\,\% at none,
40.1\,\% at 0.15\,s and 81.9\,\% at 0.3\,s dead time, 57.5\,\% at 0.2\,s latency. Right: yaw
reversals per minute on the route: 2.18 at no delay, still 2.18 at 0.15\,s dead time and
2.17 at 0.1 and 0.2\,s latency, then 7.97 at 0.25\,s latency (disturb\_v1, mean of seeds
0--2, which agree) and 9.49 at 0.3\,s dead time. The data are sparse: six delays with
route data and four with step data, one variable at a time; the lines only join them.
Dashed grey: the crane-measured $L \approx 0.165$\,s; shaded: $L > 0.3$\,s, fix latency
before tuning. Data: \code{figures/tuning/data/delay\_*.dat}.}
\label{fig:delay}
\end{figure}

\file{crane\_ml Tools/Tuning/latency.py; pwm\_relay.py (sim dead time)}
